\chapter{บทนำ}
\label{chapter:introduction}

\section{ที่มาและความสำคัญ}

ในระดับมหาวิทยาลัย ระบบการเรียนการสอนมีความซับซ้อนมากกว่าระดับมัธยมศึกษาอย่างมีนัยสำคัญ ทั้งในด้านโครงสร้างหลักสูตร การลงทะเบียนเรียน ระเบียบข้อบังคับทางวิชาการ การคำนวณผลการเรียน และขั้นตอนทางวิชาการต่าง ๆ ความซับซ้อนดังกล่าวส่งผลให้นักศึกษา โดยเฉพาะนักศึกษาใหม่ ประสบความยากลำบากในการทำความเข้าใจระบบการเรียนและการปรับตัวในช่วงเริ่มต้นของการศึกษาในระดับมหาวิทยาลัย

นอกจากนี้ ข้อมูลสำคัญทางการศึกษาของมหาวิทยาลัยมักกระจัดกระจายอยู่ในหลายแหล่ง เช่น เว็บไซต์สำนักทะเบียน ประกาศข่าวสาร เอกสารในรูปแบบไฟล์ PDF และเอกสารจากคณะหรือภาควิชา ทำให้การเข้าถึงข้อมูลเป็นไปได้ยาก ข้อมูลบางส่วนอาจไม่เป็นปัจจุบันหรือสูญหาย ส่งผลให้นักศึกษาเกิดความสับสนและอาจตัดสินใจผิดพลาด เช่น การลงทะเบียนเรียนไม่เป็นไปตามลำดับของหลักสูตร หรือการอ้างอิงระเบียบการศึกษาที่ไม่ตรงกับปีการศึกษาที่ตนสังกัด

ในด้านข้อมูลวิชาการ รายละเอียดเกี่ยวกับผังหลักสูตร รายวิชาเลือกเสรี และรายวิชาบังคับก่อนเรียน (Prerequisite) มักถูกเขียนในลักษณะเป็นทางการและมีความซับซ้อน ทำให้นักศึกษาบางส่วนตีความข้อมูลผิดพลาด และเลือกลงทะเบียนในรายวิชาที่ไม่สามารถลงเรียนได้จริง ส่งผลกระทบต่อแผนการเรียนและระยะเวลาการศึกษา

ขณะเดียวกัน ปัญหาเกี่ยวกับการคำนวณผลการเรียนและการประเมินความเสี่ยงต่อการพ้นสภาพการเป็นนักศึกษายังคงเป็นประเด็นสำคัญ นักศึกษาจำนวนมากไม่เข้าใจเกณฑ์คะแนนเฉลี่ยสะสมขั้นต่ำ (GPA) หรือเงื่อนไขการพ้นสภาพ ทำให้ไม่สามารถประเมินสถานะทางการเรียนของตนเองได้อย่างถูกต้อง ส่งผลให้ขาดการวางแผนแก้ไขปัญหาล่วงหน้า

สำหรับอาจารย์ที่ปรึกษา การติดตามและดูแลนักศึกษาก็เป็นไปอย่างจำกัด เนื่องจากข้อมูลที่เกี่ยวข้องกับผลการเรียนและความเสี่ยงของนักศึกษากระจายอยู่ในหลายระบบ และขาดเครื่องมือที่สามารถสรุปภาพรวมของนักศึกษาได้อย่างมีประสิทธิภาพ ส่งผลให้การให้คำปรึกษามักเกิดขึ้นหลังจากนักศึกษาเผชิญกับปัญหาแล้ว

จากปัญหาดังกล่าว จึงมีความจำเป็นในการพัฒนา \textbf{UniAssist AI} ที่สามารถรวบรวมข้อมูลทางการศึกษาของมหาวิทยาลัยไว้ในศูนย์กลางเดียว อธิบายข้อมูลที่มีความซับซ้อนให้อยู่ในรูปแบบที่เข้าใจง่าย ให้คำแนะนำด้านวิชาการอย่างถูกต้อง ประเมินความเสี่ยงทางการเรียนของนักศึกษาแบบเรียลไทม์ และสนับสนุนการทำงานของอาจารย์ที่ปรึกษาในการติดตามและดูแลนักศึกษาได้อย่างมีประสิทธิภาพมากยิ่งขึ้น

ในการพัฒนาระบบตอบคำถามนั้น เนื่องจากข้อมูลหลักสูตรของคณะมีลักษณะเป็นข้อมูลที่มีโครงสร้างชัดเจน (เช่น รายวิชา หน่วยกิต ปี/ภาคการศึกษา และเงื่อนไขบังคับก่อน) โครงงานนี้จึงเลือกจัดเก็บข้อมูลในรูปแบบฐานข้อมูลเชิงสัมพันธ์ (Relational Database) และประยุกต์ใช้เทคนิค \textbf{Text-to-SQL} ซึ่งให้โมเดลภาษาขนาดใหญ่ (Large Language Model: LLM) แปลงคำถามภาษาไทยของผู้ใช้ให้เป็นคำสั่ง SQL เพื่อค้นคืนข้อมูลจากฐานข้อมูลโดยตรง แนวทางนี้ทำให้คำตอบอ้างอิงจากข้อมูลจริงในฐานข้อมูลของคณะ ลดปัญหาการสร้างข้อมูลเท็จ (Hallucination) และสามารถตรวจสอบย้อนกลับได้ว่าคำตอบมาจากคำสั่งค้นข้อมูลใด

\section{วัตถุประสงค์}

\begin{enumerate}
    \item เพื่อพัฒนาระบบผู้ช่วยด้านการศึกษา (AI Education Assistant) ที่รวบรวมข้อมูลทางการศึกษาของคณะเทคโนโลยีสารสนเทศจากหลายแหล่งให้อยู่ในฐานข้อมูลเดียว เพื่อให้นักศึกษาเข้าถึงข้อมูลที่ถูกต้องและเป็นปัจจุบันได้อย่างสะดวก
    \item เพื่อพัฒนา AI Chatbot ด้านวิชาการที่ตอบคำถามของนักศึกษาได้ตลอด 24 ชั่วโมง โดยประยุกต์ใช้เทคนิค Text-to-SQL ร่วมกับความสามารถในการให้เหตุผลของโมเดลภาษาขนาดใหญ่ (LLM Reasoning) เพื่อเพิ่มความแม่นยำและความเข้าใจง่ายของคำตอบ พร้อมทั้งลดปัญหาการให้ข้อมูลเท็จ
    \item เพื่อออกแบบระบบคำนวณผลการเรียนและประเมินความเสี่ยงทางการศึกษาโดยอัตโนมัติ เช่น ความเสี่ยงต่อการพ้นสภาพการเป็นนักศึกษา พร้อมทั้งให้คำแนะนำเชิงปฏิบัติแก่นักศึกษาได้อย่างทันท่วงที
    \item เพื่อพัฒนาระบบข้อมูลทุนการศึกษาและข้อมูลด้านการเงิน พร้อมช่วยวิเคราะห์คุณสมบัติเบื้องต้นของนักศึกษาเทียบกับเงื่อนไขการสมัครทุน
    \item เพื่อสนับสนุนการทำงานของอาจารย์ที่ปรึกษา โดยจัดให้มีแดชบอร์ดติดตามผลการเรียนและการประเมินความเสี่ยงของนักศึกษาในความดูแล เพื่อเพิ่มประสิทธิภาพในการให้คำปรึกษาทางวิชาการ
    \item เพื่อส่งเสริมเป้าหมายการพัฒนาที่ยั่งยืน SDG 4: คุณภาพการศึกษา โดยเพิ่มความเท่าเทียมในการเข้าถึงข้อมูลทางการศึกษา ส่งเสริมการวางแผนการเรียนอย่างเป็นระบบ และช่วยลดปัญหาการเรียนล่าช้าในระดับอุดมศึกษา
\end{enumerate}

\section{วิธีการดำเนินงาน}

\begin{enumerate}
    \item \textbf{การตั้งปัญหาและวิเคราะห์ความต้องการของระบบ} วิเคราะห์ปัญหาและความต้องการของนักศึกษาและอาจารย์ที่ปรึกษา เช่น การเข้าถึงข้อมูลการศึกษา การวางแผนการเรียน ทุนการศึกษา และการติดตามความเสี่ยง เพื่อกำหนดขอบเขตและฟังก์ชันของระบบ
    \item \textbf{การรวบรวมและจัดเตรียมข้อมูล} รวบรวมข้อมูลทางการศึกษาจากเอกสารหลักสูตร (มคอ.2) ผังหลักสูตร ระเบียบการศึกษา และตารางสอนจากสำนักทะเบียน จากนั้นใช้โมเดลภาษาขนาดใหญ่สกัดข้อมูลที่มีโครงสร้างซับซ้อนให้อยู่ในรูปแบบข้อมูลที่มีโครงสร้าง (Structured Data) และจัดเก็บลงในฐานข้อมูลเชิงสัมพันธ์
    \item \textbf{การประเมินคุณภาพการสกัดข้อมูล} เปรียบเทียบผลการสกัดข้อมูลจากโมเดลภาษาหลายตัว (Qwen, Gemini, DeepSeek และ Typhoon) กับเอกสารหลักสูตรฉบับจริง เพื่อคัดเลือกโมเดลที่เหมาะสมที่สุดสำหรับสร้างฐานข้อมูลและใช้งานในระบบ
    \item \textbf{การพัฒนาแกน Text-to-SQL Chatbot} พัฒนาแกนการทำงานที่ให้ LLM แปลงคำถามภาษาไทยเป็นคำสั่ง SQL รันบนฐานข้อมูลแบบอ่านอย่างเดียว ผ่านกลไกตรวจสอบความปลอดภัยของคำสั่ง แล้วเรียบเรียงผลลัพธ์กลับเป็นภาษาที่เข้าใจง่าย
    \item \textbf{การพัฒนาระบบคำนวณเกรดและประเมินความเสี่ยง} พัฒนาระบบคำนวณผลการเรียนและประเมินความเสี่ยงของนักศึกษาโดยอัตโนมัติ โดยอ้างอิงกฎเกณฑ์ของสถาบันที่จัดเก็บในฐานข้อมูล และใช้ LLM สรุปสถานะและให้คำแนะนำเชิงปฏิบัติ
    \item \textbf{การพัฒนาระบบทุนการศึกษาและแดชบอร์ดอาจารย์} พัฒนาระบบให้ข้อมูลทุนการศึกษาและวิเคราะห์คุณสมบัติเบื้องต้น และพัฒนาแดชบอร์ดสำหรับอาจารย์ที่ปรึกษาเพื่อแสดงภาพรวมผลการเรียนและความเสี่ยงของนักศึกษา
    \item \textbf{การพัฒนาส่วนติดต่อผู้ใช้และระบบยืนยันตัวตน} พัฒนาเว็บแอปพลิเคชันและระบบยืนยันตัวตนผ่าน Google OAuth พร้อมแบ่งสิทธิ์การใช้งานตามบทบาท (นักศึกษา / อาจารย์ที่ปรึกษา / ผู้พัฒนา)
    \item \textbf{การพัฒนาและทดสอบระบบ} พัฒนาระบบในรูปแบบเว็บแอปพลิเคชันและดำเนินการทดสอบด้านความถูกต้องของข้อมูล การคำนวณเกรด และการตอบคำถามของ AI Chatbot เพื่อปรับปรุงระบบให้มีความน่าเชื่อถือ
\end{enumerate}

\section{ขอบเขตของโครงงาน}

\begin{enumerate}
    \item พัฒนาและรวบรวมข้อมูลด้านการศึกษาที่เกี่ยวข้องกับคณะเทคโนโลยีสารสนเทศทุกหลักสูตรในระดับปริญญาตรี (IT, DSBA, BIT และ AIT) ของสถาบันเทคโนโลยีพระจอมเกล้าเจ้าคุณทหารลาดกระบัง โดยดึงข้อมูลจากแหล่งข้อมูลที่เกี่ยวข้อง เช่น เอกสารหลักสูตร มคอ.2 ผังหลักสูตร ระเบียบการศึกษา และตารางสอนของสำนักทะเบียน
    \item พัฒนาระบบ AI Chatbot ด้านการศึกษาสำหรับนักศึกษาคณะเทคโนโลยีสารสนเทศ เพื่อตอบคำถามด้านการเรียน การลงทะเบียน ผังหลักสูตร และระเบียบการศึกษา โดยประยุกต์ใช้เทคนิค Text-to-SQL ร่วมกับความสามารถในการให้เหตุผลของโมเดลภาษาขนาดใหญ่
    \item ออกแบบและพัฒนาระบบคำนวณผลการเรียนและประเมินความเสี่ยงทางการศึกษาของนักศึกษา เช่น ความเสี่ยงต่อการพ้นสภาพการเป็นนักศึกษา โดยอ้างอิงตามเกณฑ์และข้อบังคับของสถาบัน
    \item พัฒนาระบบข้อมูลทุนการศึกษา เพื่อให้ข้อมูลเกี่ยวกับเงื่อนไขและคุณสมบัติการสมัครทุน พร้อมช่วยวิเคราะห์คุณสมบัติเบื้องต้นของนักศึกษา
    \item ออกแบบและพัฒนาแดชบอร์ดสำหรับอาจารย์ที่ปรึกษา เพื่อติดตามข้อมูลผลการเรียน วิเคราะห์ความเสี่ยง และสนับสนุนการให้คำแนะนำทางการศึกษาแก่นักศึกษาในความดูแล
    \item ออกแบบและพัฒนาส่วนติดต่อผู้ใช้ (UX/UI) ของเว็บแอปพลิเคชัน พร้อมระบบยืนยันตัวตนและการควบคุมสิทธิ์การเข้าถึงตามบทบาทผู้ใช้
    \item พัฒนาและทดสอบระบบในรูปแบบเว็บแอปพลิเคชัน โดยทดสอบความถูกต้องของข้อมูล การคำนวณผลการเรียน และการตอบคำถามของ AI Chatbot เพื่อประเมินประสิทธิภาพและปรับปรุงระบบให้เหมาะสมกับการใช้งานจริง
\end{enumerate}

\section{ประโยชน์ที่คาดว่าจะได้รับ}

\begin{enumerate}
    \item นักศึกษาสามารถติดตามข่าวสารและข้อมูลด้านการศึกษาได้อย่างรวดเร็ว และสอบถามข้อสงสัยต่าง ๆ ได้ทันทีผ่านระบบเดียว
    \item อาจารย์ที่ปรึกษาสามารถติดตามและดูแลนักศึกษาได้อย่างมีประสิทธิภาพ และรับรู้สถานะการเรียนและความเสี่ยงของนักศึกษาได้
    \item นักศึกษาสามารถวางแผนการเรียนและเลือกลงทะเบียนรายวิชาได้อย่างราบรื่นและสะดวกมากยิ่งขึ้น
    \item AI Chatbot ที่พัฒนาขึ้นสามารถตอบคำถามด้านการศึกษาได้อย่างถูกต้อง ไม่สร้างข้อมูลเท็จ และอ้างอิงจากข้อมูลจริงของคณะโดยตรง
    \item ระบบช่วยส่งเสริมการเข้าถึงข้อมูลด้านการศึกษาอย่างง่ายดาย เท่าเทียม และทั่วถึง สอดคล้องกับเป้าหมายการพัฒนาที่ยั่งยืน SDG 4 ด้านคุณภาพการศึกษา
\end{enumerate}

\section{ศัพท์นิยามเฉพาะ}

\begin{enumerate}
    \item \textbf{Large Language Model (LLM)} โมเดลภาษาขนาดใหญ่ที่ผ่านการฝึกฝนด้วยข้อมูลจำนวนมหาศาล มีความสามารถในการเข้าใจ ตีความ และสร้างข้อความภาษามนุษย์ รวมถึงมีความสามารถในการให้เหตุผล (Reasoning)
    \item \textbf{Text-to-SQL} เทคนิคการแปลงคำถามภาษาธรรมชาติ (Natural Language) ให้เป็นคำสั่งภาษา SQL โดยอัตโนมัติ เพื่อค้นคืนข้อมูลจากฐานข้อมูลเชิงสัมพันธ์ตามความต้องการของผู้ใช้
    \item \textbf{Relational Database} ฐานข้อมูลเชิงสัมพันธ์ที่จัดเก็บข้อมูลในรูปแบบตารางที่มีความสัมพันธ์กัน และเข้าถึงด้วยภาษา SQL ในโครงงานนี้ใช้ระบบจัดการฐานข้อมูล SQLite
    \item \textbf{Retrieval-Augmented Generation (RAG)} เทคนิคที่เพิ่มประสิทธิภาพให้แก่โมเดลภาษาโดยการค้นคืนข้อมูลที่เกี่ยวข้องจากฐานความรู้ภายนอกมาประกอบการสร้างคำตอบ เพื่อลดการให้ข้อมูลเท็จและทำให้ข้อมูลเป็นปัจจุบัน
    \item \textbf{Prerequisite (รายวิชาบังคับก่อน)} รายวิชาพื้นฐานหรือเงื่อนไขที่นักศึกษาจำเป็นต้องลงทะเบียนเรียนและสอบผ่านตามเกณฑ์ที่กำหนด ก่อนที่จะมีสิทธิ์ลงทะเบียนในรายวิชาต่อเนื่องหรือรายวิชาในระดับที่สูงขึ้นตามผังหลักสูตร
    \item \textbf{Academic Dismissal (การพ้นสภาพการเป็นนักศึกษา)} สถานะทางการศึกษาที่สิ้นสุดลงเนื่องจากนักศึกษามีผลการเรียนเฉลี่ยสะสม (GPA) ต่ำกว่าเกณฑ์ขั้นต่ำที่สถาบันกำหนด หรือไม่ปฏิบัติตามเงื่อนไขการศึกษาระหว่างที่เป็นนักศึกษา
    \item \textbf{Hallucination} ปรากฏการณ์ที่โมเดลภาษาสร้างคำตอบที่ฟังดูสมเหตุสมผลแต่ไม่สอดคล้องกับข้อเท็จจริงหรือไม่มีแหล่งอ้างอิงจริง
    \item \textbf{OAuth} มาตรฐานการมอบสิทธิ์ (Authorization) ที่อนุญาตให้แอปพลิเคชันยืนยันตัวตนผู้ใช้ผ่านผู้ให้บริการภายนอก (เช่น Google) โดยไม่ต้องจัดเก็บรหัสผ่านของผู้ใช้เอง
    \item \textbf{มคอ.2} เอกสารรายละเอียดของหลักสูตร (Program Specification) ตามกรอบมาตรฐานคุณวุฒิระดับอุดมศึกษาแห่งชาติ (TQF) ซึ่งระบุโครงสร้างหลักสูตร รายวิชา หน่วยกิต และผลการเรียนรู้ที่คาดหวัง
    \item \textbf{SDG 4 (Sustainable Development Goal 4)} เป้าหมายการพัฒนาที่ยั่งยืนข้อที่ 4 ของสหประชาชาติ ว่าด้วยการสร้างหลักประกันว่าทุกคนมีการศึกษาที่มีคุณภาพอย่างครอบคลุมและเท่าเทียม
\end{enumerate}
